\chapter{Introduction}
\label{ch:introduction}

The field of Information Retrieval (IR) has experienced a paradigm shift with the advent of Large Language Models (LLMs) and transformer-based architectures like BERT \cite{devlin2018bert}. These models have demonstrated unprecedented capabilities in understanding semantic relationships between user queries and documents, significantly outperforming traditional lexical ranking algorithms such as BM25. However, this increase in retrieval accuracy comes at a substantial computational cost. Fine-tuning models with hundreds of millions, or even billions, of parameters on massive datasets like MS MARCO \cite{nogueira2020document} requires immense GPU resources and memory, making it inaccessible for many researchers and practitioners.

To mitigate this computational bottleneck, Parameter-Efficient Fine-Tuning (PEFT) techniques have been developed. Methods such as Low-Rank Adaptation (LoRA) \cite{hu2021lora} have gained widespread adoption by freezing the vast majority of the network's pre-trained weights and training only a small number of injected, low-rank matrices within specific layers (typically attention computation modules). While highly effective, LoRA still requires modifications deep within the model architecture and introduces additional computational overhead during the forward pass of training.

This thesis investigates a highly constrained, embedding-level parameter-efficient fine-tuning approach as an alternative to LoRA for text ranking tasks. Specifically, we explore a "gradient masking" technique applied to a pre-trained BERT cross-encoder. Rather than injecting adapter layers or tuning entire blocks of transformer parameters, we strictly freeze the entire network architecture and isolate training entirely to the structural embeddings that guide the model's ranking decisions: the \texttt{[CLS]} (classification) and \texttt{[SEP]} (separator) tokens. 

By observing how the model adjusts solely its representation of these task-critical tokens when exposed to relevance labels, we aim to uncover the underlying learning mechanisms of sequence-to-sequence ranking described in earlier works like MonoT5 \cite{nogueira2020document}.

The primary objective of this research is to evaluate whether tuning only two token embeddings (amounting to a microscopic fraction of the total model parameters) can yield competitive ranking performance compared to a standard LoRA implementation and a fully fine-tuned baseline like MonoBERT \cite{lin2021pretrained}. Performance is robustly evaluated using the BEIR benchmark \cite{thakur2021beir}, focusing on standard metrics such as Normalized Discounted Cumulative Gain (NDCG@10) and Mean Reciprocal Rank (MRR@10).

The remainder of this thesis is structured as follows. Chapter \ref{ch:background} provides the necessary background on Information Retrieval, evaluation metrics, and Transformer architectures. Chapter \ref{ch:methodology} details the implementation of both the prompt embedding baseline and the LoRA approach. Chapter \ref{ch:experiments} describes the experimental setup, datasets, and training configurations. Chapter \ref{ch:results} presents the comparative evaluation of the models. Finally, Chapter \ref{ch:conclusions} summarises the findings and discusses avenues for future work.
